\section{Comparación, decisiones y cierre}

\begin{frame}{Comparación entre análisis horizontal y vertical}
  \begin{columns}[T,onlytextwidth]
    \column{0.48\textwidth}
    \scriptsize
    \begin{tabular}{lll}
      \toprule
      Criterio & Horizontal & Vertical \\
      \midrule
      Enfoque & Tiempo & Estructura \\
      Pregunta & ¿Cuánto cambió? & ¿Qué peso tiene? \\
      Base & Periodo anterior & Base común \\
      Utilidad & Tendencia & Composición \\
      \bottomrule
    \end{tabular}
    \tiempo{0.8 minutos}
    \column{0.48\textwidth}
    \figpng[0.62\textheight]{F12_Comparacion_Horizontal_Vertical.png}
  \end{columns}
  \note{Comparar ambos métodos.}
\end{frame}

\begin{frame}{Ejemplo comparativo en servicios digitales}
  \begin{columns}[T,onlytextwidth]
    \column{0.45\textwidth}
    \footnotesize
    Una empresa de servicios digitales observa que su gasto en servidores aumentó entre 2024 y 2025.

    \vspace{0.45em}
    El análisis horizontal muestra un aumento de S/ 4,000.00 y 50.00 \%. El análisis vertical muestra que el gasto representa 30.00 \% de las ventas.

    \vspace{0.45em}
    Ambos resultados sugieren revisar sostenibilidad y eficiencia.
    \tiempo{0.8 minutos}
    \column{0.51\textwidth}
    \figpng[0.66\textheight]{F20_Ejemplo_Comparativo_Horizontal_Vertical.png}
  \end{columns}
  \note{Ejemplo comparativo integrado.}
\end{frame}

\begin{frame}{Errores frecuentes en la interpretación}
  \begin{columns}[T,onlytextwidth]
    \column{0.46\textwidth}
    \begin{itemize}
      \item Confundir ventas con rentabilidad.
      \item Ignorar la base porcentual.
      \item Leer cifras sin contexto.
      \item No complementar con ratios.
    \end{itemize}
    \tiempo{0.8 minutos}
    \column{0.50\textwidth}
    \figpng[0.62\textheight]{F13_Errores_Frecuentes_Interpretacion.png}
  \end{columns}
  \note{Errores frecuentes.}
\end{frame}

\begin{frame}{Aplicación empresarial}
  \begin{columns}[T,onlytextwidth]
    \column{0.48\textwidth}
    \begin{itemize}
      \item Precios y costos.
      \item Financiamiento e inversión.
      \item Presupuesto.
      \item Control de gestión.
    \end{itemize}
    \micro{Los resultados del análisis permiten priorizar acciones correctivas y asignar recursos con mayor criterio.}
    \tiempo{0.8 minutos}
    \column{0.48\textwidth}
    \figpng[0.62\textheight]{F14_Decisiones_Empresariales.png}
  \end{columns}
  \note{Aplicación empresarial.}
\end{frame}

\begin{frame}{Aplicación en Ingeniería de Sistemas y análisis de datos}
  \begin{columns}[T,onlytextwidth]
    \column{0.44\textwidth}
    \begin{itemize}
      \item ERP y bases de datos.
      \item ETL y cálculo financiero.
      \item SQL, Excel, Power BI y Python.
      \item Dashboard y alerta gerencial.
    \end{itemize}
    \micro{El vínculo con Contabilidad Empresarial se mantiene porque los indicadores nacen de estados financieros y cuentas contables.}
    \tiempo{0.9 minutos}
    \column{0.52\textwidth}
    \figpng[0.66\textheight]{F17_Flujo_ERP_Dashboard.png}
  \end{columns}
  \note{Aplicación en Ingeniería de Sistemas.}
\end{frame}

\begin{frame}{Conclusiones finales}
  \footnotesize
  \begin{itemize}
    \item Los estados financieros organizan información económica y financiera en reportes comparables.
    \item El análisis horizontal identifica tendencias y variaciones entre periodos.
    \item El análisis vertical revela la estructura interna y el peso relativo de las partidas.
    \item La Empresa ABC muestra que vender más no siempre significa mayor eficiencia.
    \item La interpretación financiera orienta decisiones sobre precios, costos, inversión y financiamiento.
    \item La Ingeniería de Sistemas permite automatizar reportes, dashboards y alertas gerenciales.
  \end{itemize}
  \tiempo{1 minuto}
  \note{Conclusiones por eje.}
\end{frame}
